\section{The real place}\label{sec:real}

In this section we prove that $\Delta_N$ is small not only at $G$ but on an exponentially small interval
around it. The statement is Theorem~\ref{thm:real}: for all large $N$ and $|x-G|\le e^{-\rho N}$,
\[
 |\Delta_N(x)|\le\exp\bigl(-(n_0^2(U-\tfrac12\log2)-\varepsilon)N^2\bigr),\qquad \rho=n_0c_2+w .
\]
Here again the design is a parameter. We write $\alpha_A=A/n=(b_0+q_0+2g_0)/n_0$,
$\beta_h=2h_0/n_0$, $\gamma_H=H/n=(n_0+b_0+q_0+g_0+h_0)/n_0$, $\gamma_d=d_0/n_0=(n_0-g_0)/n_0$,
and
\[
 c_{07}=1-\frac{b_0+q_0+g_0}{n_0},\qquad c_{26}=3-\frac{b_0+q_0+2g_0}{n_0}.
\]
For $D$ these are $\alpha_A=11/29$, $\beta_h=3/29$, $\gamma_H=39/29$, $\gamma_d=53/58$, $c_{07}=41/58$ and
$c_{26}=76/29$. The proof has three parts. We first reduce the perturbation $\Delta_N(x)-\Delta_N(G)$ to
determinants in which some columns of $\cA_N$ are replaced by values of the row functions at fixed
points (\S\ref{ssec:quadrature}--\ref{ssec:andreief}). We then bound such determinants by OpenAI's
energy method, extended to configurations with fewer free points (\S\ref{ssec:energy}). Finally we
supply the trial sequences that the energy method needs and verify the resulting inequalities
(\S\ref{ssec:trials}--\ref{ssec:cells}).

\subsection{A quadrature formula for the coefficient of the constant}\label{ssec:quadrature}
Let $\delta=x-G$. Then $\Delta_N(x)=\det(\cA_N+\delta\cL_N)$, and the question is how large the
columns of $\cL_N$ are compared with those of $\cA_N$. By \eqref{eq:L} the entries of $\cL_N$ are linear
functionals of the coefficients of the row polynomials. Since these polynomials have bounded degree,
the functionals can be written exactly as finite combinations of point evaluations: if
$\tau_1,\ldots,\tau_J$ are distinct points and $J$ is at least the number of degrees of freedom of the
rows, Lagrange interpolation gives
\begin{equation}\label{eq:quad}
 \cL_N(r,k)=\sum_{a=1}^JA_r(\tau_a)\,\beta_{a,k},
\end{equation}
where the weights $\beta_{a,k}$ are the values of the same functionals on the Lagrange basis
polynomials. This is a Gauss-type quadrature formula, exact for polynomials of the relevant degree.
The point of the formula is that the columns of $\cL_N$ become combinations of point evaluations of the
same row functions that define $\cA_N$, so that the perturbed determinant can again be treated by the
energy method. The cost is the size of the weights.

We take $J=d_0N=53N$ nodes. Put $g_*=401/2500$ and $a_*=2g_*/(1+g_*^2)=2005000/6410801$, and let
$K=\lceil J/2\rceil$. The nodes are the points
\[
 t_{j,\pm}=\pm\sqrt{a_*^2+(1-a_*^2)(1+s_j)/2},\qquad T_K(s_j)=0,
\]
that is, the roots of $T_K\bigl((2t^2-(1+a_*^2))/(1-a_*^2)\bigr)$; we keep all $K$ positive roots and
the first $J-K$ negative ones, so that for odd $J$ exactly one root is dropped. They lie in
$(-1,-a_*)\cup(a_*,1)$. In the coordinate $x$ with $t=2x/(1+x^2)$ every node has $|x|>g_*$; this matters
below, because the energy fields are larger near $x=0$. Dropping a root changes the nodal polynomial by a
linear factor, and the derivative bound, the circle bound and the degree of the Chebyshev polynomial by
constant factors only, so an odd number of nodes costs nothing in the exponent. This is needed here,
since $d_0=53$ is odd and $J=53N$ is odd for odd $N$.

\begin{lemma}\label{lem:quadrature}
After removal of the common factor of the rows, the weights in \eqref{eq:quad} satisfy, for every
$w>\Wcap$ and all large $N$,
\[
 |\beta_{a,k}|\le108000N\,e^{wN}.
\]
\end{lemma}

\begin{proof}
At a root, the derivative identity for Chebyshev polynomials bounds $|P_J'(t_{j,\pm})|$ from below by a
constant multiple of $N a_*/(1-a_*^2)$, up to the dropped factor. The weight is a coefficient
functional of the Lagrange polynomial, and every coefficient of a polynomial is bounded by its maximum
modulus on a circle centred at the origin. On the circle $|z|=r$ with $r=1001/1000$, write $s=|\operatorname{Re}z|$.
For $\omega=(2z^2-(1+a_*^2))/(1-a_*^2)$ one has $|T_K(\omega)|\le(\Theta+\sqrt{\Theta^2-1})^K$, where
\[
 \Theta(s)=\frac{\sqrt{(r^2+1)^2-4s^2}+\sqrt{(r^2+a_*^2)^2-4a_*^2s^2}}{1-a_*^2}
\]
is half the sum of the distances from $\omega$ to $\pm1$. The function $\Theta$ decreases in $s$, while the
factor $|1-z|^{q_0+h_0}\le(r^2+1+2s)^{N(q_0+h_0)/2}$ coming from the column filter and from the common
factor of the rows increases. A cover of $[0,r]$ by $42$ intervals, with rational upper bounds for the
square roots at their endpoints, proves, with $d_0=53$ and $q_0+h_0=6$,
\[
 \frac{d_0}2\log\bigl(\Theta(s)+\sqrt{\Theta(s)^2-1}\bigr)+\frac{q_0+h_0}2\log(r^2+1+2s)<\Wcap\qquad(0\le s\le r),
\]
the supremum being $52.7849$ at $s\approx0.125$. The cover is done on the square of this inequality
raised to the power $N$, so that no logarithm is evaluated and the odd value of $d_0N$ costs only a
factor. Since $|z-t_{j,\pm}|\ge1/1000$ on the circle, the lemma follows.
\end{proof}

The circle of radius $1.001$ instead of $1$ separates the circle from the nodes; on the unit circle itself
the cancellation in the factor $(1-t)^h$ near $t=1$ is what keeps the bound small.

\subsection{Column expansion and an identity of Andr\'eief type}\label{ssec:andreief}
Expanding $\det(\cA_N+\delta\Phi\beta)$, where $\Phi(r,a)=A_r(\tau_a)$, multilinearly in the columns gives
a sum over maps $f$ from the columns to $\{\varnothing,1,\ldots,J\}$ of
$\prod_{f(k)\ne\varnothing}\delta\beta_{f(k),k}\det M_f$. Here $M_f$ is the matrix whose $k$-th column is the
$k$-th column of $\cA_N$ if $f(k)=\varnothing$ and the vector $(A_r(\tau_{f(k)}))_r$ otherwise.

\begin{lemma}\label{lem:expansion}
If $|\beta_{a,k}|\le B$ and $|\det M_f|\le Mc^{m(f)}$ for every $f$, where $m(f)$ is the number of columns
replaced, then
\[
 \bigl|\det(\cA_N+\delta\Phi\beta)-\det\cA_N\bigr|\le M\bigl((1+J|\delta|Bc)^{n}-1\bigr).
\]
\end{lemma}

Each $M_f$ has $p=n-m$ columns that are integrals and $m$ columns that are point values. Andr\'eief's
identity extends to this situation.

\begin{lemma}\label{lem:andreief}
Let $\mu,\nu$ be finite measures, $F_r$ $(r<p+m)$, $G_i$ $(i<p)$ and $K$ integrable functions, and let $B$
be the matrix with $B(r,i)=\iint F_r(t)K(t,s)G_i(s)\,d\mu(t)\,d\nu(s)$ for $i<p$ and
$B(r,p+j)=F_r(\tau'_j)$. Then
\begin{multline*}
 (p!)^2\det B=\int\!\!\int\det\bigl[F_r((t\Vert\tau')_c)\bigr]_{r,c}\\
 \cdot\det[K(t_i,s_j)]_{i,j<p}\det[G_i(s_j)]_{j,i<p}\,d\mu^p(t)\,d\nu^p(s),
\end{multline*}
where $t\Vert\tau'$ is the list of the $p$ variables followed by the $m$ fixed points.
\end{lemma}

\begin{proof}
For $m=0$ this is Andr\'eief's identity applied twice \cite{Forrester2018Andreief}. In general expand
$\det B$ multilinearly in its first $p$ columns, which introduces one pair of variables $(t_i,s_i)$ for
each of them, and antisymmetrize first in the $t$ and then in the $s$ variables; the last $m$ columns are
constant throughout.
\end{proof}

With $K(t,s)=1/(1-ts)$ the middle determinant is the Cauchy determinant, and the third is a
generalized Vandermonde determinant $\prod_js_j^b(1-s_j)^q\det[s_j^{E_i}]$. For the latter we use the
bound
\begin{equation}\label{eq:genvand}
 \bigl|\det[s_j^{E_i}]\bigr|\le\Bigl(\prod_j(1+s_j)\Bigr)^{n-p}\prod_{i<j}|s_i-s_j|,
\end{equation}
valid for distinct exponents $E_i<n$ and $s_j\in[0,1]$. It follows by multiplying the matrix by a
unitriangular matrix built from the polynomial $\prod_j(X-s_j)$ and bounding the remaining $(n-p)\times(n-p)$
determinant by the product of the $\ell^1$ norms of its columns. The quotient
$\det[s_j^{E_i}]/\prod_{i<j}(s_j-s_i)$ is a Schur polynomial, so the bound is close to optimal.

\subsection{The energy estimate with fixed points}\label{ssec:energy}
OpenAI bounds the integrand in Andr\'eief's identity by the exponential of $n^2E$, where $E$ is a discrete
logarithmic energy of the points $x_i$ and $s_j$ in the coordinate $t=2x/(1+x^2)$. The energy consists of
an external field for each point, logarithmic interactions $\log|x_i-x_j|$, $\log(1-x_ix_j)$ and
$\log|s_i-s_j|$ among points of the same kind, and an interaction $-\log(1-2x_is_j+x_i^2)$ between the two
kinds. There are two cases, according to which of OpenAI's two bounds for the determinant of the row
functions is used \cite[Props.~5.2, 5.3]{OpenAI2026}; in the second case the interaction among the $x_i$
has weight $2$ instead of $1$, and in the first the row field acquires a term with a multiplier
$\lambda_1$ for the constraint that defines the case.

The energy is bounded from above by duality. Given trial sequences $u=(u_k)$ and $v=(v_k)$, one writes
the interaction terms as quadratic forms in the moments $\frac1n\sum_iT_k(x_i)$ and
$\frac1n\sum_js_j^k$ and completes the squares against $u$ and $v$. What remains is the sum of the norms
$\|u\|^2=\sum_ku_k^2/k$, $\|v\|^2$ and, for each point, a function of one variable. With
$\mathsf T(u,x)=\sum_ku_kT_k(x)/k$ and $\mathsf S(v,x)=\sum_kv_kx^k/k$ these are, in the second case,
\begin{align*}
 X(x)&=\alpha_A\log|x|+\beta_h\log(1-x)-\gamma_H\log(1+x^2)-4\mathsf T(u,x)-\mathsf S(v,x),\\
 Y(s)&=\tfrac{b_0}{n_0}\log s+\tfrac{q_0}{n_0}\log(1-s)+2\mathsf T(v,s),
\end{align*}
and in the first case
\begin{multline*}
 X_1(x)=\alpha_A\log|x|+\beta_h\log(1-x)-(\gamma_H-\tfrac12)\log(1+x^2)\\
 +\lambda_1\bigl(\tfrac{2g_0}{n_0}-\tfrac{2x^2}{1+x^2}\bigr)
 -2\mathsf T(p,x)-\mathsf S(v,x),
\end{multline*}
with the same $Y$ (for the first case the trial sequence for the first variables is called $p$).
If $X\le M_X$ and $Y\le M_Y$ on the whole domains, the energy of any configuration is at most a
constant determined by $M_X$, $M_Y$ and the norms. In potential-theoretic terms the extremal
configuration is described by an equilibrium measure for two interacting measures in an external
field \cite{SaffTotik}, and the best trial sequences are its moments, as Carreira observed
\cite{Carreira2026}; any trial sequences give a valid upper bound.

When $m$ of the $n$ columns are point values, the configuration has only $n-m$ free points of each kind,
while the normalization is still by $n$. The fixed points contribute through the determinant of the row
functions evaluated at the concatenated list $t\Vert\tau'$. The fixed points are of two kinds: $a$ of
them are quadrature nodes, with $|x|>g_*$, and $b'$ of them are points $y\in(0,1)$; the numbers of free
points, of the first kind and of the second kind add up to $n$, and we put $\alpha=(a+b')/n$ and
$\tau=b'/n$. The missing free points change
the energy bound by an amount linear in $m$, which is the cost of a fixed column. Every step of OpenAI's
argument remains valid uniformly in the number of free points: the regularization of the logarithmic
interactions, the diagonal corrections, and the completion of squares only involve sums over the points
present. For the fixed points of the first kind the relevant field is the residual field
\[
 R(x)=-\gamma_d\log(1+x^2)-2\mathsf T(u,x),\qquad g_*\le|x|\le1,
\]
and for those of the second kind it is
\[
 N(y)=-\gamma_d\log(1+y^2)+\frac{n_0-2g_0-1}{n_0}\log y-u_1y,\qquad0<y<1,
\]
where $u_1$ is the first coefficient of $u$; the coefficient $(n_0-2g_0-1)/n_0$ is a lower bound, valid for
all $N\ge1$, of $(n-1-2g)/n$.

\begin{lemma}[budget, second case]\label{lem:budget2}
Let $\ell=\log2$, and let $\mathcal N,\Phi,\mathcal R,\mathcal N_c$ and $c_2$ be real numbers, with $m\ge0$,
such that
\[
 \mathcal N+\Phi+U-c_{07}\ell\le-m,\qquad(c_{26}\ell-\Phi+\mathcal R-c_2)^2\le8\ell m,\qquad
 \ell-\mathcal R+\mathcal N_c\le0 .
\]
Then for all $\alpha,\tau\ge0$
\begin{multline*}
 \bigl(-(c_{07}-\tfrac12)+(c_{26}-1)\alpha-\alpha^2+\tau-\tau^2\bigr)\ell+\mathcal N+(1-\alpha)\Phi+\alpha(1-\alpha)\ell\\
 +(\alpha-\tau)\mathcal R+\tau\mathcal N_c\le-U+\tfrac12\ell+c_2\alpha .
\end{multline*}
\end{lemma}

\begin{proof}
Subtracting the right side from the left gives
$(\mathcal N+\Phi+U-c_{07}\ell)+\alpha(c_{26}\ell-\Phi+\mathcal R-c_2)-2\ell\alpha^2
+\tau(\ell-\mathcal R+\mathcal N_c)-\ell\tau^2$. The last two terms are non-positive by the third hypothesis,
and the first three form a concave quadratic in $\alpha$ with non-positive discriminant, since
$A+s\alpha-2\ell\alpha^2\le0$ for all $\alpha$ as soon as $s^2+8\ell A\le0$.
\end{proof}

The lemma is applied with $\mathcal N$ a bound for $2\|u\|^2+\frac12\|v\|^2$, $\Phi=M_X+M_Y$, $\mathcal R$ a
bound for the supremum of $R$ and $\mathcal N_c$ a bound for the supremum of $N$. The analogue for the first
case has no fixed points of the second kind, and its polynomial in $\alpha$ is
\begin{multline*}
 \bigl(-(c_{07}-\tfrac12)+(c_{26}-1)\alpha-\alpha^2\bigr)\ell+n_1+(1-\alpha)F_1+\alpha(1-\alpha)\ell\\
 +\alpha\lambda_1\tfrac{2g_0}{n_0}\le-U+\tfrac12\ell+c_1\alpha,
\end{multline*}
with $n_1$ a bound for the norms of the first trial and $F_1=M_{X_1}+M_Y$; it holds for all real
$\alpha$ when its discriminant is non-positive and $c_1\le c_{26}\ell-F_1+\lambda_1\tfrac{2g_0}{n_0}$ (so that
it also holds for $\alpha<0$).

Combining the energy bounds with Lemma~\ref{lem:quadrature}, a determinant $M_f$ with $m$ fixed columns
contributes at most $(5\cdot108000N)^me^{(n_0c+w)Nm}$ times the unperturbed bound, where $c=\max(c_1,c_2)$;
the factor $5^m$ comes from expanding the row functions at the fixed points in the second case.
Lemma~\ref{lem:expansion} then gives, for $|\delta|\le e^{-\rho N}$ with $\rho=n_0c_2+w$, a factor
$(1+O(N^2))^{n}=e^{o(N^2)}$ (here $w$ is any number above $\Wcap$ and $c=c_2$). The unperturbed bound is
$\exp\bigl(n_0^2N^2(-U+\frac12\log2+\varepsilon)\bigr)$. This proves Theorem~\ref{thm:real}, given the
numerical bounds of the next two subsections.

\subsection{The trial sequences}\label{ssec:trials}
For the first case we use a trial sequence with $40$ Chebyshev coefficients and a finite number of
geometric terms, with the multiplier $\lambda_1=1.68493357$. For the second case we use trial sequences
of the form
\begin{equation}\label{eq:trial}
 u_k=c_k+\operatorname{Re}\sum_jr_jz_j^k\qquad(k\ge1;\ c_k=0\text{ for }k>K),
\end{equation}
with $K=\Kcoef$ Chebyshev coefficients, dyadic rationals, and $\Ngeo$ geometric terms, $\Ncomplex$ of them with
complex $z_j$ and complex $r_j$.

The geometric terms have a simple meaning. Since $\sum_{k\ge1}z^kT_k(x)/k=-\frac12\log(1-2xz+z^2)$, a term
$(z,r)$ contributes $-\frac12\operatorname{Re}\,r\log(1-2xz+z^2)$ to the potential $\mathsf T(u,x)$. If
$x_0+i\epsilon=\frac12(z+z^{-1})$ then $1-2xz+z^2=2z(x_0+i\epsilon-x)$, so the term is, up to a constant, the
logarithmic potential of a point charge of complex strength $r$ at $x_0+i\epsilon$ and its conjugate. The
extremal measure of the energy problem has a gap in its support around $x=0$, and its potential has a
logarithmic singularity at $x=0$ and square-root behaviour at the edges of the gap. A Chebyshev expansion
converges slowly near such points. A charge at a small distance $\epsilon$ from them reproduces the local
behaviour with one term. We place the complex charges at
$x_0\in\{0,\pm0.03,\pm0.06,\pm0.09,\pm0.12,\pm0.16\}$ with $\epsilon\in\{10^{-2},3\cdot10^{-3}\}$, and
the real terms at points $\pm z$ with $z$ increasing from $0.8$ to $0.999997$; the numbers $z_j$ are
rationals with at most twelve decimals. The norms of the geometric terms have closed forms: two terms $(z_j,r_j)$,
$(z_l,r_l)$ contribute
$-\frac12\operatorname{Re}\bigl(r_jr_l\log(1-z_jz_l)+r_j\bar r_l\log(1-z_j\bar z_l)\bigr)$ to $\|u\|^2$.

\subsection{Verification of the bounds}\label{ssec:cells}
We must bound three functions of one variable from above on their whole domains: the field $X$ on
$[-1,1]$, the field $Y$ on $[0,1]$, and the residual field $R$ on $g_*\le|x|\le1$. With $x=\cos\theta$,
each is a constant plus a cosine series with $K$ terms plus finitely many terms
$\operatorname{Re}\,w\log(1-ae^{im\theta})$.

We cover each domain by cells $[x_c-h,x_c+h]$ whose endpoints are multiples of $2^{-40}$, and on
each cell bound the function by its Taylor expansion at $x_c$. The cells are grouped in blocks. The
polynomial part has degree $K-1$, and a global bound for its fifth derivative would force very short cells.
Instead, for each block the derivatives of the polynomial part are computed exactly at the centre of the
block up to order $70$, and the remainder beyond order $70$ is bounded by a single constant for the block,
which enters multiplied by $h^{70}/70!$; the Taylor expansion at the centre of a cell is obtained from
that of the block by an exact shift. The logarithmic terms are expanded to order $5$, with remainders
bounded from the exact minimum of $|1-ae^{im\theta}|$ on the cell. A few cells at $x=0$ and $x=\pm1$ treat
the singular terms separately. The covers of $X$, $Y$ and $R$ have \Nblocks{} blocks and \Ncells{} cells, respectively.
Each cell is checked by exact rational and integer arithmetic.

The resulting bounds are
\[
 \sup X\le\SupX,\quad \sup Y\le\SupY,\quad
 \sup_{g_*\le|x|\le1}R\le\ResCap,
\]
so that $M_X+M_Y=\FieldCap$; the norms satisfy $2\|u\|^2+\frac12\|v\|^2\le\NormCap$; and the first
coefficient satisfies $u_1\ge9/50$, which is used near $x=0$. For the near points, a one-variable
estimate shows that
$\sup_{0<y<1}N(y)\le-77/100$ for $u_1\ge9/50$, with the weight $47/58$ in $N(y)$; the true supremum is
about $-0.7758$. For the first case, with the weights $(\alpha_A,\beta_h,\gamma_H-\frac12,\lambda_1,2g_0/n_0)=(11/29,3/29,49/58,\lambda_1,5/29)$,
the bounds $\sup X_1\le-1.15854$ and $\sup Y\le-1.32981$ are obtained by isolating the zeros of the
numerators of the derivatives, which are polynomials with integer coefficients, and bounding the values
between them; the norm bound is $n_1=0.67761$.

With $c_{07}=41/58$, $c_{26}=76/29$ and $U=\Utarget$, the hypotheses of Lemma~\ref{lem:budget2} hold with
\begin{gather*}
 m=\Mmargin,\qquad \ell-\mathcal R+\mathcal N_c\le-0.0137,\\
 (c_{26}\ell-\Phi+\mathcal R-c_2)^2\le\Disc\le8\ell m=\DiscBound,
\end{gather*}
with $c_2=\Ctwo$; for the first case $c_1=\Cone<c_2$ is admissible. This gives the cost $c_2$ of a fixed
column, and $\rho=58c_2+w=\Rho$ for $w=2.9\cdot18.203=52.7887$ ($>\Wcap$). The decay rate in
Theorem~\ref{thm:real} is $n_0^2(U-\frac12\log2)=\TDelta\ldots$ The trial data are fixed once and for
all; the file containing them has SHA-256 digest
\begin{center}
{\footnotesize\texttt{\Digest}}.
\end{center}
